% ID: SEL-D0004-OA-E04
% ESTADO: cotejada. Cotejo literal independiente y revisión de composición realizados.
% Fuente: Selectividad/sel_2019_matematicas_aplicadas/Suplente Septiembre/Exámenes Matemáticas Aplicadas CCSS - A y B.pdf
% SHA-256: eecad753694b3d8cb2d427efddb00d75faae9c4ec5b027056291bc95d7d2e3da
% Página: 1; opción A; ejercicio 4.
% Encabezado original: EJERCICIO 4. Solo enunciado; no contiene solución.

A la salida de una heladería se realizó una encuesta para comprobar si los
clientes habían probado un nuevo sabor en promoción. Se observó que de 125
personas encuestadas, 20 no lo habían probado y el resto sí.

\begin{enumerate}[label=\alph*)]
  \item \textbf{(1.5 puntos)} Determine, con un nivel de confianza del
  97\,\%, un intervalo para estimar la proporción de clientes de esa
  heladería que no habían probado el nuevo helado.
  \item \textbf{(1 punto)} Mediante una nueva muestra se desea estimar la
  proporción de clientes de esa heladería que no habían probado el nuevo
  helado, con un error inferior al 5\,\% y un nivel de confianza del
  94\,\%. Suponiendo que se mantiene la proporción muestral del apartado
  anterior, \textquestiondown{}qué tamaño mínimo debe tener dicha muestra?
\end{enumerate}
